\documentclass[10pt]{article}
\usepackage[T1]{fontenc}
\usepackage[utf8]{inputenc}
\usepackage[french]{babel}
\usepackage{amsmath,amssymb,mathrsfs,mathpazo}
\usepackage{geometry}
\geometry{a4paper,margin=2cm}
\usepackage{graphicx,array,tabularx,tabularray,colortbl,multirow,multicol,booktabs,enumitem,pifont,ulem,fancybox,diagbox,cancel}
\usepackage{pstricks,pstricks-add,pst-plot,pst-tree,pst-node,pst-text,pst-func,pst-math,pst-bspline,pst-3dplot,pst-eucl,pst-solides3d}
\let\pscalculate\undefined
\usepackage{pst-bezier}
\usepackage{tikz,pgfplots,tkz-tab}
\usetikzlibrary{arrows,arrows.meta,patterns,calc,positioning,trees,decorations.markings,shapes.geometric,shapes.misc,shapes.arrows,intersections,decorations.pathmorphing,decorations.pathreplacing,backgrounds,shadows,babel}
\pgfplotsset{compat=1.18}
\usepgfplotslibrary{fillbetween}
\usepackage[np]{numprint}
\usepackage{listings,makecell,pas-tableur,twoopt}
\let\barre\undefined
\usepackage[pstricks]{tablvar}
\definecolor{darkgreen}{RGB}{0,100,0}
\providecommand{\red}{\color{red}}
\providecommand{\blue}{\color{blue}}
\providecommand{\green}{\color{green}}
\providecommand{\black}{\color{black}}
\providecommand{\white}{\color{white}}
\providecommand{\gray}{\color{gray}}
\providecommand{\cyan}{\color{cyan}}
\providecommand{\magenta}{\color{magenta}}
\providecommand{\yellow}{\color{yellow}}
\providecommand{\violet}{\color{violet}}
\providecommand{\orange}{\color{orange}}
\providecommand{\euro}{\texteuro}
\providecommand{\EUR}{\texteuro}
\providecommand{\ieme}{\textsuperscript{e}}
\providecommand{\degres}{\ensuremath{^\circ}}
\providecommand{\degree}{\ensuremath{^\circ}}
\providecommand{\vv}[1]{\overrightarrow{#1}}
\providecommand{\Coord}[1]{#1}
\providecommand{\decosix}{}
\providecommand{\starredbullet}{\ensuremath{\bullet}}
\pagestyle{empty}
\setlength{\parindent}{0pt}
\usepackage{adjustbox}\DeclareUnicodeCharacter{2236}{:}\DeclareUnicodeCharacter{2190}{\ensuremath{\leftarrow}}
\providecommand{\R}{\mathbb{R}}
\providecommand{\N}{\mathbb{N}}
\providecommand{\Z}{\mathbb{Z}}
\providecommand{\Q}{\mathbb{Q}}
\providecommand{\D}{\mathbb{D}}
\providecommand{\C}{\mathbb{C}}
\providecommand{\e}{{\mathrm{e}}}
\providecommand{\vect}[1]{\overrightarrow{#1}}
\providecommand{\vectt}[1]{\overrightarrow{\mathrm{#1}}}
\providecommand{\vvt}[1]{\overrightarrow{\mathrm{#1}}}
\providecommand{\np}[2][]{{#2~#1}}
\providecommand{\vv}[1]{{\overrightarrow{#1}}}
\renewcommand{\d}{{\mathrm{d}}}
\providecommand{\degres}{{^\circ}}
\providecommand{\degre}{{^\circ}}
\providecommand{\blue}{\color{blue}}
\providecommand{\red}{\color{red}}
\providecommand{\enskip}{\quad}
\providecommand{\llbracket}{\lbrack\!\lbrack}
\providecommand{\rrbracket}{\rbrack\!\rbrack}
\providecommand{\euro}{{\text{€}}}
\providecommand{\no}{{\text{n°}}}
\providecommand{\interval}[2]{{[#1;#2]}}
\providecommand{\intervFO}[2]{{[#1;#2[}}
\providecommand{\dotfill}{\ldots}

\providecommand{\ds}{\displaystyle}
\providecommand{\textding}[1]{#1}
\providecommand{\Oij}{\(\mathrm{O};\vec{\imath},\vec{\jmath}\)}
\providecommand{\Oijk}{\(\mathrm{O};\vec{\imath},\vec{\jmath},\vec{k}\)}
\providecommand{\Ouv}{\(\mathrm{O};\vec{u},\vec{v}\)}
\providecommand{\intervFF}[2]{\ensuremath{[#1;#2]}}
\renewcommand{\euro}{\ifmmode\text{\texteuro}\else\texteuro\fi}
\DeclareUnicodeCharacter{00B0}{\ifmmode^{\circ}\else\textdegree\fi}
\DeclareUnicodeCharacter{20AC}{\ifmmode\text{\texteuro}\else\texteuro\fi}
% Énoncés : Éducation nationale. Transcription et corrigés : APMEP, Denis Vergès et contributeurs. Découpage et classement : CoopMaths, Sébastien Lozano, Éric Elter et contributeurs. Adaptation de la présentation : Maths Explorer.
% CC BY-SA : https://coopmaths.fr/www/about/
% https://coopmaths.fr/alea/static/bac/2023/tex/bac_2023_03_sujet2_asie_4_cor.tex
\begin{document}
\noindent\textit{APMEP / CoopMaths. CC BY-SA. Présentation adaptée par Maths Explorer.}\par\medskip
\medskip

\emph{Pour chacune des cinq questions de cet exercice, une seule des quatre réponses proposées est exacte.\\Le candidat indiquera sur sa copie le numéro de la question et la réponse choisie.\\
Aucune justification n'est demandée.\\
Une réponse fausse, une réponse multiple ou l'absence de réponse à une question ne rapporte ni n'enlève de point.}

\medskip





\textbf{Question 1 :} Le nombre de termes de cette liste est : 673 : en effet,

le $n$-ième terme de la liste est $7 + 3(n - 1)$ ; il faut donc résoudre l'équation :

$\np{2023} = 7 + 3(n - 1) \iff \np{2016} = 3(n - 1) \iff 3 \times 672 = 3(n - 1) \iff 672 = n - 1 \iff n = 673$.








\textbf{Question 2 :} On choisit au hasard un nombre dans cette liste. La probabilité de tirer un nombre pair est:







Il faut trouver $n$ tel que $7 + 3(n - 1) = 2k$, avec $k \in \N$,soit $7 + 3n - 3 = 2k \iff 4 + 3n = 2k$.

Il faut que la somme $4 + 3n$ soit paire ; comme 4 est pair il faut donc que $3n$ soit pair et comme 3 est impair il faut que $n$ soit pair.

Sur les 673 nombres de la liste le premier et le dernier nombre de la liste sont impairs ; il y a donc 336 nombres pairs. Réponse C
\medskip

On rappelle qu'on choisit au hasard un nombre dans cette liste. 

On s'intéresse aux évènements suivants :
\begin{itemize}
\item[$\bullet~~$]Évènement $A$ : \og obtenir un multiple de 4 \fg
\item[$\bullet~~$]Évènement $B$ : \og obtenir un nombre dont le chiffre des unités est 6 \fg
\end{itemize}

Pour répondre aux questions suivantes on pourra utiliser l'arbre pondéré ci-dessous et on donne
$p(A \cap B) = \dfrac{34}{673}$.

\begin{center}
\pstree[treemode=R,nodesepA=0pt,nodesepB=2.5pt,treesep = 1cm,levelsep=2.5cm]{\TR{}}
{\pstree{\TR{$A$~~}\taput{\footnotesize$\dfrac{168}{673}$}}
{\TR{$B$~~}\taput{\footnotesize $\dfrac{17}{84}$}
\TR{$\overline{B}$}\tbput{\footnotesize $\dfrac{67}{84}$}
}
\pstree{\TR{$\overline{A}$~~}\tbput{\footnotesize$\dfrac{505}{673}$}}
{\TR{$B$~~}\taput{\footnotesize $\dfrac{33}{505}$}
\TR{$\overline{B}$}\tbput{\footnotesize $\dfrac{472}{505}$}
}
}
\end{center}

\medskip

\textbf{Question 3 :}

La probabilité d'obtenir un multiple de 4 ayant 6 comme chiffre des unités est donnée par 

$p(A \cap B) = \dfrac{34}{673}$ ! Réponse B







\medskip

\textbf{Question 4 :} $P_B(A)$ est égale à :

$P_B(A) = \dfrac{P(B \cap A)}{P(B)}$.

Or $P(B) = P(A \cap B) + P\left(\overline{A} \cap B \right) = \dfrac{34}{673} + \dfrac{505}{673} \times
\dfrac{33}{505} = \dfrac{34}{673} + \dfrac{33}{673}  = \dfrac{67}{673}$.

Donc $P_B(A) = \dfrac{\dfrac{34}{673}}{\dfrac{67}{673}} = \dfrac{34}{67}$. Réponse D








\medskip

\textbf{Question 5:} On choisit, au hasard, successivement, $10$ éléments de cette liste. 

Un élément peut être choisi plusieurs fois. La probabilité qu'aucun de ces $10$ nombres ne soit un multiple de 4 est:








La probabilité de ne pas tirer un multiple de 4 est égale à $\dfrac{505}{673}$.

La probabilité qu'aucun des 10 nombres tirés ne soit un multiple de 4 est égale à $\left(\dfrac{505}{673}\right)^{10} \approx 0,057$.
Réponse A

\end{document}
